\subsection{Entorno y método}

El laboratorio se ejecutó en un equipo con Windows, usando Docker Desktop
(con backend WSL2/Linux) para levantar dos contenedores (\texttt{cliente} y
\texttt{servidor}) conectados por la red \texttt{sd\_net}, siguiendo el
material provisto por el curso (\texttt{docker-compose.yml},
\texttt{cliente.py}, \texttt{servidor.py}).

Durante la puesta en marcha, Docker Desktop reportó el error
\texttt{Virtualization support not detected}. La virtualización por
hardware ya estaba habilitada en el equipo; lo que faltaba eran las
características de Windows necesarias para que Docker pudiera usarla
(\texttt{VirtualMachinePlatform} y el subsistema WSL). Se solucionó
ejecutando PowerShell como administrador y corriendo:

\begin{verbatim}
dism.exe /online /enable-feature /featurename:VirtualMachinePlatform /all /norestart
dism.exe /online /enable-feature /featurename:Microsoft-Windows-Subsystem-Linux /
all /norestart
bcdedit /set hypervisorlaunchtype auto
\end{verbatim}

seguido de un reinicio del equipo, tras lo cual Docker Desktop instaló el
soporte de WSL faltante e inició correctamente.

La medición se realizó en cuatro pasos: (1) línea base con \texttt{ping} e
\texttt{iperf3} sobre la red sin modificar (2) 100 y 1000 llamadas
PING/PONG secuenciales con la red sana (3) inyección de latencia con
\texttt{tc qdisc netem delay} en 50, 200 y 500\,ms, midiendo el total de 100
llamadas en cada nivel (4) inyección de pérdida de paquetes con
\texttt{tc qdisc netem loss} en 1\%, 5\%, 20\% y 100\%, midiendo throughput
con \texttt{iperf3} y comportamiento del cliente en cada nivel, incluyendo
la falla provocada con \texttt{loss 100\%}.

\subsection{Resultados}

\begin{table}[h]
\centering
\begin{tabular}{|l|c|c|c|}
\hline
\textbf{Latencia \texttt{tc}} & \textbf{Total (s)} & \textbf{Promedio (ms)} & \textbf{Máx (ms)} \\
\hline
0 ms (base) & 0.011 & 0.1 & 1.4 \\
50 ms  & 5.041  & 50.4  & 51.3 \\
200 ms & 20.045 & 200.4 & 200.9 \\
500 ms & 50.046 & 500.5 & 501.0 \\
\hline
\end{tabular}
\caption{Latencia inyectada vs. tiempo total de 100 llamadas}
\end{table}

\begin{table}[h]
\centering
\begin{tabular}{|l|c|c|}
\hline
\textbf{Pérdida \texttt{tc}} & \textbf{Throughput iperf3} & \textbf{Total cliente (s)} \\
\hline
1\%   & 10.7 Gbit/s & 0.011 \\
5\%   & 449 Mbit/s  & 1.459 \\
20\%  & 1.21 Mbit/s & 5.252 \\
100\% & --          & error inmediato / timeout \\
\hline
\end{tabular}
\caption{Pérdida de paquetes inyectada vs. throughput y comportamiento del cliente}
\end{table}

% Gráfico: total (s) vs latencia (ms) — relación lineal, pendiente ~0.1s por cada ms de latencia
% Insertar aquí figura o \includegraphics con el gráfico de latencia vs total

\subsection{Observado vs esperado}

En el Paso 3 (latencia), los cálculos teóricos ($N \times \text{RTT}$ para
100 llamadas secuenciales) coincidieron casi exactamente con lo medido en
los tres niveles: 50\,ms $\rightarrow$ 5.0\,s esperado vs 5.041\,s medido;
200\,ms $\rightarrow$ 20\,s esperado vs 20.045\,s medido; 500\,ms
$\rightarrow$ 50\,s esperado vs 50.046\,s medido. La degradación fue
perfectamente lineal, sin sorpresas: el cliente espera cada respuesta antes
de enviar la siguiente llamada, por lo que el retardo se acumula de forma
directa.

En el Paso 4 (pérdida) las diferencias fueron mayores a lo esperado
intuitivamente. Con solo 1\% de pérdida, el throughput de \texttt{iperf3}
ya cayó a menos de la mitad (de 23.2\,Gbit/s a 10.7\,Gbit/s), lo que no es
proporcional al porcentaje de pérdida: se explica porque TCP reduce
agresivamente su ventana de congestión ante la primera señal de pérdida.
En cambio, el cliente de 100 llamadas no mostró cambios visibles con 1\% de
pérdida (total casi igual al base), porque con pocos paquetes chicos es
probable que ninguno se perdiera en esa corrida particular. Esto muestra
que la cantidad de tráfico determina si una falla de baja probabilidad
estadística llega a manifestarse.

\subsection{La falacia que asumimos}

La falacia identificada es \textbf{``La red es confiable''}
(Deutsch/Gosling, falacia \#1). Se ve claramente en \texttt{cliente.py},
en la línea

\begin{verbatim}
TIMEOUT = float(os.getenv("TIMEOUT_S", "0")) or None
\end{verbatim}

donde, al no definirse \texttt{TIMEOUT\_S}, la expresión \texttt{0 or None}
se evalúa como \texttt{None}, dejando el socket \textbf{sin timeout} en la
llamada \texttt{socket.create\_connection}. Esto se traduce directamente en
la línea donde el cliente espera cada respuesta:

\begin{verbatim}
respuesta = f.readline()  # aquí el cliente ESPERA a la red
\end{verbatim}

\texttt{readline()} es una llamada bloqueante: si el servidor no responde,
el programa se queda esperando indefinidamente, sin ningún mecanismo
propio para detectar que la red falló. Esto se comprobó en la falla
provocada (\texttt{loss 100\%}): sin \texttt{TIMEOUT\_S}, el programa no
mostró ningún error de aplicación propio dependió de un error externo del
sistema (\texttt{Errno 113 No route to host}) con \texttt{TIMEOUT\_S=3},
el cliente abortó de forma controlada a los 3.0\,s, mostrando que el
mecanismo de timeout sí existe en el código, pero \textbf{no está activado
por defecto}.

El cambio que corregiría esta falacia es simple: cambiar el valor por
defecto de \texttt{TIMEOUT\_S} de \texttt{"0"} a un valor explícito
razonable (por ejemplo, 5 segundos), de modo que la ausencia de
configuración no implique esperar para siempre. Además, se podría agregar
lógica de reintento con backoff exponencial y un número máximo de
intentos antes de reportar la falla al usuario, en vez de terminar con un
único error tras el timeout.

\subsection{Relación con la investigación}

La aplicación investigada para la Evaluación 1 es \textbf{Transbank / Webpay}, uno de los pagos electrónicos más usada en Chile. En su flujo de integración
(Webpay Plus), el comercio debe realizar una llamada de \emph{confirmación}
(\texttt{commit}) al servidor de Transbank tras la autorización del pago para
saber si la transacción fue aprobada.

La falacia que más ha afectado a este tipo de sistema es \textbf{``la red es
confiable''}, la misma identificada en el Paso 4 de este laboratorio. Si la
llamada de confirmación falla por timeout o pérdida de paquetes, el comercio
no puede asumir que la transacción no ocurrió: el cargo puede haberse
efectuado igualmente del lado del banco emisor, y solo se perdió la
respuesta de vuelta, el mismo fenómeno de retransmisión silenciosa que se
observó con \texttt{tc netem loss} en este laboratorio, donde el cliente
completaba llamadas con éxito pero con máximos disparados sin que la
aplicación se enterara del retraso real. Por eso la documentación de
integración de Transbank recomienda no asumir un resultado negativo ante un
timeout, sino volver a consultar el estado real de la transacción mediante
un endpoint de \texttt{status} antes de decidir si reintentar el cobro.
Un manejo ingenuo de este timeout (asumiendo directamente que ``no llegó
respuesta = no se cobró'') puede llevar a cobros duplicados o a discrepancias
entre lo que el comercio y el banco registran como transacción exitosa.